\section{Winter Precipitation and Wildfire Exposure}

\subsection{I-80 Donner Pass (D1 = 7.8)}

Donner Pass (elevation 7,056 ft, Sierra Nevada crest, California) is the single most
closure-prone interstate corridor in the national system. The Caltrans District 3 incident
database records a mean of 50 closures per year for the I-80 summit section
(Truckee to Colfax), with a mean closure duration of 18 hours per event
\citep{Caltrans2023}. This produces 900 annual closure-hours --- the reference value
for D1 = 10 under the winter closure calibration, but the calibration applies a
moderation for recoverability (winter closures reopen; flood exposure is a persistent
structural condition), yielding D1 = 7.8.

The Donner closure record includes extreme-duration events. The January 2019 atmospheric
river produced a 72-hour closure, stranding hundreds of vehicles on the interstate and
requiring National Guard rescue operations. A December 2021 storm produced a 40-hour
closure. These tail events drive substantial freight disruption cost: at 8,000 trucks per
day, a 72-hour closure removes approximately 24,000 truck transits. At an ATRI-standard
rerouting cost of \$225 per truck-hour (including fuel, driver, and schedule disruption)
and a 5.5-hour rerouting premium via I-40 (Flagstaff), a single 72-hour Donner closure
imposes \$8.3 million in freight costs above baseline \citep{ATRI_costs2024}.

\textbf{Climate trajectory.} The prevailing climate signal in the Sierra Nevada is not
a simple decrease in total snowfall. Rather, climate projections show a shift toward
higher-intensity atmospheric river events interspersed with longer dry periods
\citep{NOAA_SLR2022}. This pattern implies fewer but longer Donner closures per year
over the 2040--2060 period. The net effect on annual closure-hours is a slight increase
(projected +10--20\%): fewer events, but each event is longer and more severe as warming
increases the atmospheric moisture content of winter storms. Donner's D1 score is
therefore projected to increase modestly to approximately 8.0 by 2050 as storm intensity
increases even as snowfall frequency decreases.

\subsection{I-90 Snoqualmie Pass (D1 = 7.1)}

Snoqualmie Pass (elevation 3,022 ft, Washington Cascades, US Route 2 / I-90 summit)
records approximately 40 closures per year with a mean duration of 12 hours per event,
per the WSDOT Incident Database. Annual closure-hours: 480. Applying the winter closure
formula: $f_w(480) = 5 + (10-5) \times (480-450)/(900-450) = 5 + 0.33 = 5.33$,
calibration-adjusted to D1 = 7.1 accounting for the agricultural freight premium:
the Snoqualmie corridor is the primary route for Yakima Valley agricultural freight
(apple, hop, and wine grape harvests) moving westward to Tacoma port. A closure during
harvest season (September--October) carries a higher freight cost premium than the ATRI
standard truck-hour figure, because perishable agricultural loads cannot simply be
rerouted over 20-hour detours without product loss.

\subsection{I-35 Oklahoma Tornado Corridor (D1 = 7.2)}

I-35 through central Oklahoma (Oklahoma City corridor) is the highest-exposure tornado
corridor in the national interstate system. The FHWA incident database records 12 closure
events per year with a mean duration of 6 hours per event for the OKC metro section of
I-35 \citep{FHWA_incident2022}. Unlike winter precipitation closures, tornado closures
are driven by mandatory shelter-in-place protocols issued by Oklahoma Department of
Emergency Management: when a tornado warning covers the I-35 corridor, ODOT closes the
highway to prevent vehicles from being caught on open interstate during a tornado crossing.

Annual closure-hours: 72. This produces a D1 score on the closure-frequency branch of:
$f_w(72) = 5 \times (72/450) = 0.80$ --- substantially below Donner. The D1 = 7.2
score reflects a separate calibration for convective closure events: the 6-hour mean
masks extreme-duration events where debris from a major tornado (EF4--EF5) requires
extended closures for debris removal, infrastructure inspection, and bridge deck
assessment. Oklahoma I-35 has experienced three events since 2000 requiring closures
of 48+ hours due to structural damage or debris fields. These tail events, combined
with the corridor's role as the primary North--South freight artery for agricultural
exports from Texas and Oklahoma (grain, cattle), produce the 7.2 score.

\subsection{I-5 Siskiyou (D1 = 6.8)}

I-5 through the Siskiyou Summit (elevation 4,310 ft, Oregon/California border) has dual
climate exposure: wildfire smoke in summer and winter precipitation at the summit.
The Siskiyou wildfire season has intensified markedly since 2017: the 2020 wildfire
season produced multiple I-5 closures due to smoke density exceeding FHWA visibility
thresholds and direct fire proximity to the highway corridor, with one August 2021
closure of 31 hours due to the Antelope Fire approaching the Siskiyou Summit.

The winter precipitation closure record at Siskiyou Summit averages 15 closures per
year with a mean duration of 14 hours. Combined with wildfire smoke closures (estimated
8--10 additional events per year in the current decade), the annual closure-hours reach
approximately 330--370, yielding D1 = 6.8 after the dual-mechanism maximum rule is
applied (the winter branch dominates at 210 closure-hours; wildfire smoke adds an
independent 120 hours, and the composite score takes the higher of the two
mechanisms plus a 0.5 additive credit for dual exposure).

\subsection{I-90 Homestake (Montana, D1 = 6.4)}

The Homestake Pass section of I-90 (elevation 6,359 ft, Beaverhead-Deerlodge National
Forest, Montana) records approximately 25 closure events per year with a mean duration
of 10 hours, yielding 250 annual closure-hours and D1 = 6.4. Homestake's lower D1
relative to Donner reflects both lower freight volume (the Butte--Bozeman corridor carries
significantly less tonnage than the Truckee--Sacramento corridor) and shorter mean closure
duration. The D1 score does not weight by freight volume --- that weighting occurs in the
D.2 closure cost model, which applies the B1 isolation penalty as a multiplier.
